\introduction

\subsection{Motivation}
Industrial alloys such as steels and aluminium grades are chosen for
load-bearing parts on the basis of their elastic response. The Young's
modulus $E$ sets the stiffness and the Poisson's ratio $\nu$ sets the
transverse contraction under load. For a cubic crystal both quantities
follow from just two independent stiffness constants, $C_{11}$ and
$C_{12}$. These constants size almost every structural calculation an
engineer runs: beam deflection scales as $1/E$, the buckling load and the
natural frequency scale with $E$, and the bulk modulus is
$K = E/[3(1-2\nu)]$. At present they are measured experimentally, which is
slow and expensive. A computational route that was both accurate and cheap
would let new compositions be screened before any metal is cast.

Molecular dynamics is the natural candidate for that route, but an MD
prediction is only as good as the interatomic potential behind it.

\subsection{Elastic Constants and Engineering Design}
A part fails or survives based on two numbers. Young's modulus $E=\sigma/\varepsilon$
fixes how much it deflects under load, and Poisson's ratio
$\nu=-\varepsilon_\mathrm{lat}/\varepsilon_\mathrm{axial}$ fixes how much it
contracts sideways while it does. For a cubic crystal both reduce to the two
independent stiffness constants $C_{11}$ and $C_{12}$, which is why this work
targets those constants directly rather than $E$ and $\nu$. The same pair
propagates everywhere downstream: a beam deflects as $1/E$, a column buckles
at a load proportional to $E$, a part rings at a frequency proportional to
$\sqrt{E/\rho}$, and the bulk modulus is $K=E/[3(1-2\nu)]$. There is also a
hard physical constraint that the simulations must respect. A cubic crystal is
only mechanically stable when $C_{11}>C_{12}$, the Born--Cauchy condition, and
a potential that violates it describes a solid that would spontaneously shear
apart. This single inequality reappears as the stability filter throughout the
pipeline.

Table~\ref{tab:applications} gives a sense of the range these constants span
across real structural materials, from a soft magnesium casting alloy near
$45$\,GPa to nickel and steel near $200$\,GPa. Aluminium alloys, the design
target of this project, sit near $71$\,GPa with $\nu\approx0.33$.

\begin{table}[htbp]
	\centering
	\caption{Young's modulus and Poisson's ratio for representative structural
	materials and their typical applications. Aluminium~7075, the running
	example in this report, is highlighted in the first row.}
	\label{tab:applications}
	\renewcommand{\arraystretch}{1.2}
	\begin{tabular}{|l|l|c|c|}
		\hline
		\textbf{Application} & \textbf{Material} & \textbf{$E$ (GPa)} & \textbf{$\nu$} \\ \hline
		Aerospace fuselage   & Al~7075          & 71--72 & 0.33 \\ \hline
		Turbine disk         & Inconel 718      & 200    & 0.29 \\ \hline
		Orthopaedic implant  & Ti-6Al-4V        & 110    & 0.34 \\ \hline
		Automotive casting   & Mg AZ91          & 45     & 0.35 \\ \hline
		Structural frame     & Mild steel       & 200    & 0.30 \\ \hline
		MEMS resonator       & Si (single xtal) & 130    & 0.28 \\ \hline
	\end{tabular}
\end{table}

\subsection{Interatomic Potentials and Density Functional Theory}
An interatomic potential is a closed-form approximation of the many-body
potential energy $U(\{\mathbf{r}_i\})$, from which the forces follow as
$\mathbf{F}_i=-\nabla_i U$. The forms used in practice grew more capable over
time. The earliest were simple pairwise functions like Lennard--Jones and
Morse, which are too crude for metals because they ignore the way a bond
weakens as an atom gains more neighbours. The Embedded Atom Method introduced
in the 1980s fixed that by adding an embedding energy that depends on the local
electron density, and it works well for close-packed FCC metals. The Modified
Embedded Atom Method \citep{baskes1992,lee2001} adds an angular dependence to
that density, which lets one functional form cover FCC, BCC, HCP, and the
mixed covalent-metallic bonding of silicon. That breadth, plus its analytic
and interpretable parameters, is why this project builds on MEAM rather than a
black-box alternative.

The reference data that the MEAM parameters are fit against comes from density
functional theory. DFT solves the Kohn--Sham equations \citep{kohn1965} for the
electronic ground state, here in a plane-wave plus pseudopotential basis using
Quantum Espresso \citep{giannozzi2009}. From it the project takes the cohesive
energy $E_\mathrm{coh}$, the equilibrium lattice constant $a_0$ and bulk
modulus $B$, the single-crystal elastic constants $C_{11}$ and $C_{12}$, and
the binary formation energies $E_\mathrm{form}(i,j)$ that seed the cross-terms.
The two methods are complementary. DFT is accurate but caps out near $10^3$
atoms, while MEAM-driven MD is empirical but scales past $10^6$ atoms, so DFT
supplies the small, trustworthy reference set against which the cheap potential
is calibrated.

\subsection{The Cross-Interaction Problem}
The obstacle to using MEAM on real alloys is coverage. A published parameter
set typically spans three or four elements, for example the Al--Mg--Zn
potential of \citet{Dickel_2018} or the Al--Si--Mg--Cu--Fe set of
\citet{jelinek2012}, and says nothing about how those elements interact with
species outside its scope. A practically relevant alloy such as Al~7075
contains Al, Zn, Mg, Cu, Cr, Mn, Fe, Si, and Ti, and no single published
file spans that list. Concatenating parameters from different libraries
leaves the binary cross-terms undefined, which produces unstable or
unphysical simulations.

The literature offers two ways around this. The first is to fit a fresh MEAM
set by hand against experimental and reference data \citep{jelinek2012}. That
is accurate but slow, and the work does not transfer to a new element list.
The second is to drop MEAM for a fully data-driven machine-learning potential
\citep{behler2007,shapeev2016,mishin2021}. Those are flexible, but they need
$10^{4}$ to $10^{5}$ reference configurations and they give up the analytic
readability that makes MEAM useful in an engineering setting.

\subsection{Approach Adopted in This Project}
This project takes a hybrid route. Rather than replacing MEAM, it uses a
neural network as a \emph{surrogate} that refits the MEAM parameters
themselves. Single-element parameters are anchored to DFT
\citep{kohn1965,giannozzi2009} references, and the missing cross-terms are
adjusted so the potential reproduces target elastic constants. This keeps the
efficiency and transparency of MEAM while borrowing DFT accuracy to fill the
gaps. The five-stage pipeline that implements the idea, configuration
generation, MD evaluation, DFT reference calculation, MEAM initialization,
and neural-network optimization, is described stage by stage in
Section~\ref{sec:progress}. Stages 1 through 4 seed the potential; Stage~5
closes the loop. A separate composition-only surrogate, described in
Section~\ref{sec:comp_nn}, runs alongside Stage~2 as a consistency check and
as a standalone predictor.

\subsection{Research Question}
\textit{Can DFT reference calculations combined with neural-network surrogate
optimization produce multi-element MEAM interatomic potentials that
reproduce the elastic properties of arbitrary metallic alloys, even when the
constituent elements do not share a common published MEAM
parameterization?} The results in Sections~\ref{sec:nnip}
and~\ref{sec:validation} answer this in the affirmative inside the
aluminium-rich design region, where both surrogates land within a few percent
of experiment, and identify the loss of accuracy outside that region as the
main item of remaining work.
